\begin{tikzpicture}[
    train/.style={pub/train},
    val/.style={pub/validation},
    test/.style={pub/test},
    arrow/.style={pub/arrow},
    ann/.style={pub/heading},
]

% One representative LOGO fold. The caption supplies the full figure title.
\foreach \c in {1,...,16} {
    \pgfmathsetmacro{\col}{mod(\c-1,8)}
    \pgfmathsetmacro{\row}{floor((\c-1)/8)}
    \pgfmathsetmacro{\x}{\col*1.10}
    \pgfmathsetmacro{\y}{0.45 - \row*1.12}
    \ifnum\c=1
        \node[test] (C\c) at (\x,\y) {C\c\\TEST};
    \else\ifnum\c=8
        \node[val] (C\c) at (\x,\y) {C\c\\VAL};
    \else
        \node[train] (C\c) at (\x,\y) {C\c};
    \fi\fi
}

\node[ann] (annTest) at (0.0,-2.35) {Held-out test condition};
\draw[arrow] (annTest.north) -- (C1.south);
\node[ann] (annVal) at (4.7,-2.35) {Validation condition};
\draw[arrow] (annVal.north) -- (C8.south);

\node[pub/base, fill=white, text width=7.5cm] at (3.85,-3.65) {
  The remaining 14 conditions train the model. Repeat the split 16 times so each condition is tested once.
};

\node[train, minimum width=0.8cm, minimum height=0.45cm] (legT) at (0.2,-4.8) {};
\node[right=0.08cm of legT, pub/legend] {Train (14)};
\node[val, minimum width=0.8cm, minimum height=0.45cm] (legV) at (3.0,-4.8) {};
\node[right=0.08cm of legV, pub/legend] {Validation (1)};
\node[test, minimum width=0.8cm, minimum height=0.45cm] (legTe) at (6.3,-4.8) {};
\node[right=0.08cm of legTe, pub/legend] {Test (1)};
\end{tikzpicture}%
